\documentclass[11pt]{article}
\usepackage[margin=1in]{geometry}
\usepackage{amsmath,amssymb,amsthm}
\usepackage{booktabs,tabularx,longtable,array}
\usepackage{enumitem}
\usepackage[hidelinks]{hyperref}

\newtheorem{theorem}{Theorem}
\newtheorem{lemma}[theorem]{Lemma}
\newtheorem{proposition}[theorem]{Proposition}
\newtheorem{corollary}[theorem]{Corollary}
\newtheorem{conjecture}[theorem]{Conjecture}
\theoremstyle{definition}
\newtheorem{definition}[theorem]{Definition}
\newtheorem{assumption}{Assumption}
\renewcommand{\theassumption}{A\arabic{assumption}}
\theoremstyle{remark}
\newtheorem{remark}[theorem]{Remark}

\newcommand{\lexlt}{\mathrel{<_{\mathrm{lex}}}}
\newcommand{\lexle}{\mathrel{\le_{\mathrm{lex}}}}
\newcommand{\cfs}{\mathrm{CFS}}
\newcommand{\fl}[1]{\left\lfloor #1 \right\rfloor}
\newcommand{\ce}[1]{\left\lceil #1 \right\rceil}
\newcommand{\Fu}{\mathsf{F}}
\newcommand{\rnd}[1]{\llbracket #1 \rrbracket}
\newcommand{\code}[1]{\texttt{#1}}
\newcommand{\SC}{\mathcal S}
\providecommand{\llbracket}{[\![}
\providecommand{\rrbracket}{]\!]}

\title{RateQ: Formal Model, Correctness, Instance Optimality and\\ Exact Maintenance of Cheapest-Feasible-Service Queries}
\author{Technical companion to the RateQ paper (artifact document)}
\date{}

\emergencystretch=3em
\begin{document}
\maketitle

\begin{abstract}
This document gives complete proofs for the guarantees claimed for RateQ. We
formalise rate rule programs exactly as the reference implementation evaluates them,
including minimum billable tiers, dimensional weight, rounding rules, extrapolation,
surcharges with eligibility predicates, fee groups, fuel and percentage-of-total rules.
We prove that the suffix-minimum hull bound never exceeds the true cost (Lemma~\ref{lem:L3}),
also in its surcharge-class refinement (Lemma~\ref{lem:L3plus}). We prove that the threshold
query returns exactly the brute-force answer with the smaller-identifier tie rule, and
that it is instance optimal in a precisely stated access model (Theorem~\ref{thm:T4}). We
prove that the top-2 repair over a support set yields exactly the state of a full
recomputation for cell, band, table, rule, zone and fuel edits (Theorem~\ref{thm:T7}), and we
state construction, space, query and maintenance bounds. Results on envelope complexity
and on general linear predicates are proved where possible and stated as conjectures
otherwise. Appendix~\ref{app:code} maps every definition to the code, and
Appendix~\ref{app:mismatch} lists every point where the code and the theory differ.
\end{abstract}

\section{Model}\label{sec:model}

\subsection{Parcels and features}
A parcel is $x=(\kappa,a,L,M,S,\mathit{res},d)$ with integer key
$\kappa=\mathit{origin}\cdot 10^5+\mathit{zip}$, actual weight $a\ge 0$, sorted dimensions
$L\ge M\ge S\ge 0$, residential flag $\mathit{res}$ and ship day $d\in\mathbb Z$. Its
\emph{parcel-level features} are $L,M,S$, $G=L+2M+2S$, $V=LMS$, $L+M+S$, $a$,
$\mathit{res}$, $\mathit{zip}$, $d$ and the month of $d$. Three further features depend on
the service $s$: the billable weight $T$, the zone $\mathit{ZONE}$ and the delivery-area
level $\mathit{DAS}$. A \emph{predicate} is any Boolean function of the feature vector. The
implementation uses disjunctions of conjunctions of threshold, set-membership and range
conditions. No result below requires monotonicity of predicates.

\subsection{Rate rule programs}
Services are identified by integers $s\in\SC=\{0,\dots,n-1\}$. We fix the tolerance
$\epsilon=10^{-9}$ that the implementation uses for weight comparisons.

\begin{definition}[Rate rule program]\label{def:program}
A service $s$ consists of the following components.
\begin{enumerate}[label=(\roman*),leftmargin=*]
\item \emph{Zone map.} Starts $\zeta_{s,1}<\dots<\zeta_{s,R_s}$ and labels
$\ell_{s,i}\in\{0,1,\dots\}\cup\{\bot\}$. Then $Z_s(\kappa)=\ell_{s,i}$ for
$\kappa\in[\zeta_{s,i},\zeta_{s,i+1})$, with $\zeta_{s,R_s+1}=+\infty$, and $Z_s(\kappa)=\bot$ for
$\kappa<\zeta_{s,1}$. The intervals are closed on the left and open on the right, and $\bot$
means that the key is not served.
\item \emph{Delivery-area map} $\mathit{DAS}_s(\kappa)$, defined in the same way, with value $0$
below the first start.
\item \emph{Group key} $\gamma(s)=(D_s,\theta_s,\rho_s,\phi_s,q_s,\mathit{carry}_s,\mathit{mode}_s)$:
divisor, dimensional threshold, rounding rule, fuel class and the rounding grid, carry
and weight mode of the enterprise rule language.
\item \emph{Pre-tier billable weight} $t_s(x)=\tau_{\gamma(s)}(a,V)$, where
\begin{itemize}[leftmargin=*]
\item $\rho=\text{carrier}$: $\tau=\max(a,[V>\theta]\,V/D)$;
\item $\rho=\text{erp-style}$: $\tau=\max(\bar a,\bar v)$, where $\bar a$ rounds $a$ down to
$0.1$ and then up to an integer, and $\bar v$ is the analogous rounding of
$\lceil 1000V/D\rceil/1000$ if this exceeds $1$, the value $1$ otherwise, and $0$ unless
$V>\theta$;
\item $\rho=\text{erp}$: $\tau=r_{q,\mathit{carry}}(y)$ with $y=\max(a,v)$, $a$ or $v$ for modes
$1,2,3$, where $v=V/D$ if $D>0$ and $v=0$ otherwise, and $r_{q,\mathit{carry}}$ rounds up or
half up to the grid $q$.
\end{itemize}
\item \emph{Minimum billable tiers} $(\pi^{(2)}_s,m^{(2)}_s)$ (lower tier) and
$(\pi^{(1)}_s,m^{(1)}_s)$, each optional. Starting from $t=t_s(x)$, the tiers are applied in
this order as $t\leftarrow\max(t,m^{(j)}_s)$ if $\pi^{(j)}_s$ holds on the feature vector with
$T=t$. The result is the \emph{effective weight} $\hat t_s(x)$.
\item \emph{Rate table of upper-closed threshold pieces.} Breakpoints
$b_{s,1}<\dots<b_{s,K_s}$ define the bands $(b_{s,k-1},b_{s,k}]$, with $b_{s,0}=-\infty$, and
prices $P_s[z,k]\in\mathbb N\cup\{\infty\}$ for every zone $z$ and band $k$. The band of a
weight $t$ is
\[
k_s(t)=\min\bigl(\{k:b_{s,k}\ge t-\epsilon\}\cup\{K_s\}\bigr).
\]
\item \emph{Extrapolation.} Either none, or a limit $\lambda_s$ above which the base price
is extrapolated linearly from the top band.
\item \emph{Ineligibility predicate} $\psi_s$, which may be absent.
\item \emph{Surcharge rules} $r\in\mathcal R_s$, each with a predicate $\pi_r$, a kind, a key
$(g_r,\varphi_r)$ consisting of a fee group $g_r$ and a fuel flag $\varphi_r\in\{0,1\}$,
per-zone amounts $A_r[z]$ and a coefficient $\mathit{coef}_r$. The kinds are \emph{fixed}
(amount $A_r[z]$), \emph{per weight} (amount $\rnd{\mathit{coef}_r\,\hat t_s}$),
\emph{percentage of base} (amount $\rnd{\mathit{coef}_r B_s/100}$) and \emph{percentage of
total} (coefficient $\mathit{coef}_r$). Here $\rnd{\cdot}$ rounds to the nearest integer, with
ties to even.
\end{enumerate}
\end{definition}

The fuel schedule of class $\phi$ consists of starts $\sigma_{\phi,0}<\dots<\sigma_{\phi,W_\phi-1}$
and ratios $f_{\phi,w}$ in basis points. The ratio on day $d$ is
$f_\phi(d)=f_{\phi,\max(0,w(d))}$, where $w(d)=\max\{w:\sigma_{\phi,w}\le d\}$ and $w(d)=-1$ if
no start precedes $d$. Thus the first segment extends to $-\infty$. A class without a
schedule has $f\equiv 0$.

\begin{definition}[Cost]\label{def:cost}
Fix $x$ and $s$, and let $z=Z_s(\kappa)$, $\hat t=\hat t_s(x)$ and $f=f_{\phi_s}(d)$. All
predicates of rules and $\psi_s$ are evaluated on the feature vector with $T=\hat t$,
$\mathit{ZONE}=z$ and $\mathit{DAS}=\mathit{DAS}_s(\kappa)$. The \emph{base price} is
\[
B_s(x)=\begin{cases}
\ce{\min(P_s[z,K_s]\,\hat t/b_{s,K_s},\,10^{15})-10^{-6}} & \text{if $s$ extrapolates and } \hat t>\lambda_s,\\
P_s[z,k_s(\hat t)] & \text{otherwise.}
\end{cases}
\]
In the first case $B_s=\infty$ if $P_s[z,K_s]=\infty$. Service $s$ is \emph{eligible} for $x$
iff $z\ne\bot$, $\psi_s(x)$ is false, $B_s(x)<\infty$, and either $s$ extrapolates or
$\hat t\le b_{s,K_s}$. For a key $(g,\varphi)$ let
$\mu_{g,\varphi}=\max\bigl(\{0\}\cup\{\alpha_r: r\text{ fires},\ (g_r,\varphi_r)=(g,\varphi),\ r\text{ not
percentage of total}\}\bigr)$, where $\alpha_r$ is the amount of $r$. Let
$\Sigma^\Fu=\sum_{g}\mu_{g,1}$, $\Sigma^N=\sum_g\mu_{g,0}$, and let $Q$ be the sum of the
coefficients of the firing percentage-of-total rules. Then
\begin{equation}\label{eq:cost}
c^0_s(x)=\fl{\frac{(B_s+\Sigma^\Fu)(10^4+f)+5000}{10^4}}+\Sigma^N,\qquad
c_s(x)=c^0_s(x)+\rnd{c^0_s(x)\,Q/100},
\end{equation}
if $s$ is eligible, and $c_s(x)=\infty$ otherwise.
\end{definition}

\begin{definition}[CFS query]\label{def:cfs}
Order pairs lexicographically: $(c,s)\lexlt(c',s')$ iff $c<c'$, or $c=c'$ and $s<s'$. Let
$F(x)=\{(c_s(x),s):c_s(x)<\infty\}$. Then $\cfs(x)=\min_{\mathrm{lex}}F(x)$ if
$F(x)\ne\emptyset$, and $\cfs(x)=\bot$ otherwise. The implementation encodes $\bot$ as
$(\infty,-1)$.
\end{definition}

\subsection{Assumptions}
Every result states which of the following assumptions it uses. Appendix~\ref{app:code}
reports the check of each assumption against the instances.

\begin{assumption}[Nonnegativity]\label{as:nonneg}
Finite prices and fixed amounts are nonnegative. The coefficients of all surcharge kinds
and all fuel ratios are nonnegative. Moreover $a,V\ge 0$.
\end{assumption}
\begin{assumption}[Magnitude]\label{as:mag}
Finite prices are at most $10^{15}$, and every finite cost is smaller than the sentinel
$2^{60}$ that encodes $\infty$.
\end{assumption}
\begin{assumption}[Extrapolation limit]\label{as:extrap}
If $s$ extrapolates, then $\lambda_s\ge b_{s,K_s}$.
\end{assumption}
\begin{assumption}[Identifiers]\label{as:sid}
The service list is ordered by identifier, and the service at position $i$ has
identifier $i$.
\end{assumption}
\begin{assumption}[Tier predicates]\label{as:tier}
Tier predicates read only parcel-level features, that is, not $T$, $\mathit{ZONE}$ or
$\mathit{DAS}$. No tier minimum $m^{(j)}_s$ lies in $(b,b+\epsilon]$ for a breakpoint $b$ of
$s$.
\end{assumption}
\begin{assumption}[Breakpoint spacing]\label{as:spacing}
Distinct elements of $E_\gamma$ (Definition~\ref{def:elem}) differ by more than $\epsilon$.
\end{assumption}
\begin{assumption}[Shared delivery areas]\label{as:das}
If the surcharge-class index treats $\mathit{DAS}$ as a parcel-level feature for a group,
then all services of the group have the same delivery-area map.
\end{assumption}

Assumptions \ref{as:nonneg}--\ref{as:sid} are used for the query results. Assumptions
\ref{as:tier}--\ref{as:spacing} are used only for maintenance, and Assumption~\ref{as:das}
only for the surcharge-class index.

\section{The Index}\label{sec:index}

\begin{definition}[Key partition and signatures]\label{def:sig}
Let $C=\{0\}\cup\bigcup_s\{\zeta_{s,i}\}$ with elements $u_1<\dots<u_N$. The elementary key
intervals are $[u_i,u_{i+1})$. The \emph{signature} of interval $i$ is the vector
$\sigma_i=(Z_s(u_i))_{s\in\SC}$, and $\Sigma$ is the set of distinct signatures. For a group
$\gamma$, the \emph{local signature} is the restriction $\sigma|_\gamma$, and $\Sigma_\gamma$
is the set of distinct local signatures.
\end{definition}

\begin{proposition}[Signature partition, P1]\label{prop:P1}
Every $Z_s$ is constant on every elementary key interval, and
$|\Sigma|\le N\le 1+\sum_sR_s$. Point location takes $O(\log N)$ time, and the bound on $N$
is attained when all starts are distinct and positive.
\end{proposition}
\begin{proof}
Every start of every zone map is a cut, so no zone map changes inside an elementary
interval. Keys below $u_1=0$ do not occur. The cuts number at most $1+\sum_sR_s$, and
$\Sigma$ is a set of interval labels. Binary search over the cuts locates an interval. If
all starts are distinct and positive, all $1+\sum_sR_s$ cuts are distinct.
\end{proof}

\begin{lemma}[Group invariance, L2]\label{lem:L2}
For all services $s,s'$ in the same group, $t_s(x)=t_{s'}(x)$ and
$f_{\phi_s}(d)=f_{\phi_{s'}}(d)$ for every parcel. Hence the order of any list of services
of one group by $(h,s)$ (Definition~\ref{def:hull}) does not depend on the fuel schedule,
and a fuel edit changes no list of the index.
\end{lemma}
\begin{proof}
The function $\tau_\gamma$ of Definition~\ref{def:program}(iv) reads only the parameters in
$\gamma$, and the fuel ratio reads only $\phi\in\gamma$. The lists store $h$, which does
not involve $f$.
\end{proof}

\begin{definition}[Elementary weight intervals]\label{def:elem}
For a group $\gamma$, let $E_\gamma=\bigcup_{s\in\gamma}\{b_{s,1},\dots,b_{s,K_s}\}$ with elements
$E_1<\dots<E_m$, and let $E_0=-\infty$ and $E_{m+1}=+\infty$. The elementary interval of a
weight $t$ is $e_\gamma(t)=\min\{e\in\{1,\dots,m+1\}:E_e\ge t-\epsilon\}$.
\end{definition}

\begin{lemma}[Band--interval correspondence]\label{lem:band}
Let $s\in\gamma$, let $t$ be a weight and let $e=e_\gamma(t)$. Then:
\begin{enumerate}[label=(\alph*)]
\item $\{k:b_{s,k}\ge t-\epsilon\}=\{k:b_{s,k}\ge E_e\}$. In particular, if this set is nonempty,
then $k_s(t)=\min\{k:b_{s,k}\ge E_e\}$ and $b_{s,k_s(t)-1}<E_e\le b_{s,k_s(t)}$.
\item $j_s(e):=\min\{k:b_{s,k}\ge E_e-\epsilon\}$ satisfies $j_s(e)\le k_s(t)$ whenever the set in
(a) is nonempty.
\item $e_\gamma$ and $k_s$ are nondecreasing, and $e_\gamma(\max(t,t'))=\max(e_\gamma(t),e_\gamma(t'))$.
\item If the set in (a) is empty, then $t-\epsilon>b_{s,K_s}$.
\end{enumerate}
\end{lemma}
\begin{proof}
(a) If $b_{s,k}\ge t-\epsilon$, then $b_{s,k}$ is an element of $E_\gamma$ that is at least
$t-\epsilon$; since $E_e$ is the least such element, $b_{s,k}\ge E_e$. Conversely,
$b_{s,k}\ge E_e\ge t-\epsilon$. For $e=m+1$ both sets are empty. The second claim follows
because $k_s(t)$ is the least element of the set.
(b) $b_{s,k}\ge E_e$ implies $b_{s,k}\ge E_e-\epsilon$.
(c) Both functions are minima of sets that shrink as $t$ grows, and the identity follows
from monotonicity.
(d) If no breakpoint is at least $t-\epsilon$, then in particular $b_{s,K_s}<t-\epsilon$.
\end{proof}

\begin{definition}[Hull bound and lists]\label{def:hull}
For $s\in\gamma$, a zone $z$ and $e\in\{1,\dots,m+1\}$, let
\[
h_s(z,e)=\begin{cases}
\infty & z=\bot,\\
\min_{k\ge j_s(e)}P_s[z,k] & e\le m \text{ and } j_s(e) \text{ is defined},\\
P_s[z,K_s] & \text{otherwise, if $s$ extrapolates},\\
\infty & \text{otherwise}.
\end{cases}
\]
For $\sigma\in\Sigma_\gamma$ and $e$, the list $\mathcal L_\gamma(\sigma,e)$ contains all services of
$\gamma$ sorted by $(h_s(\sigma_s,e),s)$. For a parcel $x$, let
$z_s=Z_s(\kappa)$, $e_x=e_\gamma(t_\gamma(x))$ and $h_s(x)=h_s(z_s,e_x)$, and define the
fuel-scaled bound
\[
\beta_s(x)=\fl{\frac{h_s(x)\,(10^4+f)+5000}{10^4}},\qquad \beta_s(x)=\infty\text{ if }h_s(x)=\infty.
\]
\end{definition}

\begin{remark}
For an extrapolating service and $e\le m$ with $j_s(e)$ defined, the implementation also
takes the minimum with $P_s[z,K_s]$. This has no effect, because $K_s\ge j_s(e)$.
\end{remark}

\section{Lower Bound Correctness (L3)}\label{sec:L3}

\begin{lemma}[Base bound]\label{lem:base}
Assume \ref{as:nonneg}--\ref{as:extrap}. Let $s\in\gamma$ be eligible for $x$, and let $t''$ be
any weight with $t''\le\hat t_s(x)$. Then $B_s(x)\ge h_s(z_s,e_\gamma(t''))$.
\end{lemma}
\begin{proof}
Write $\hat t=\hat t_s(x)$, $e=e_\gamma(t'')$, $z=z_s\ne\bot$, $P=P_s[z,\cdot]$ and $K=K_s$. We
first note a fact about extrapolation. If $s$ extrapolates and $\hat t>\lambda_s$, then
$\hat t>b_{s,K}$ by \ref{as:extrap}, so $P[K]\hat t/b_{s,K}\ge P[K]$. Since $P[K]\le 10^{15}$
by \ref{as:mag}, the minimum with $10^{15}$ is at least $P[K]$. Since $P[K]$ is an integer,
$\lceil y-10^{-6}\rceil\ge P[K]$ for every $y\ge P[K]$. Hence
\begin{equation}\label{eq:ext}
B_s(x)\ge P[K]\qquad\text{whenever $s$ extrapolates and $\hat t>\lambda_s$.}
\end{equation}
\emph{Case 1: $e\le m$ and $j=j_s(e)$ is defined.} Then $h=\min_{k\ge j}P[k]$. If $s$
extrapolates and $\hat t>\lambda_s$, then (\ref{eq:ext}) and $K\ge j$ give $B_s\ge P[K]\ge h$.
Otherwise $B_s=P[k_s(\hat t)]$. If $\{k:b_{s,k}\ge t''-\epsilon\}$ is nonempty, then
$j\le k_s(t'')\le k_s(\hat t)$ by Lemma~\ref{lem:band}(b),(c). If it is empty, then
$k_s(\hat t)=K\ge j$, because $t''\le\hat t$ and the set for $\hat t$ is empty as well. In both
cases $B_s=P[k_s(\hat t)]\ge\min_{k\ge j}P[k]=h$.

\emph{Case 2: $e=m+1$, or $j_s(e)$ is undefined.} We claim $\hat t-\epsilon>b_{s,K}$. If
$e=m+1$, then $E_m<t''-\epsilon$ and $b_{s,K}\le E_m$. If $e\le m$ and $j_s(e)$ is undefined,
then $b_{s,K}<E_e-\epsilon<E_e$. Since $b_{s,K}\in E_\gamma$, we get $b_{s,K}\le E_{e-1}<t''-\epsilon$
by the definition of $e$. In both cases $\hat t-\epsilon\ge t''-\epsilon>b_{s,K}$. Hence
$\hat t>b_{s,K}$, and a non-extrapolating $s$ would be ineligible. So $s$ extrapolates and
$h=P[K]$. If $\hat t>\lambda_s$, then (\ref{eq:ext}) applies. Otherwise no breakpoint is at
least $\hat t-\epsilon$, so $k_s(\hat t)=K$ and $B_s=P[K]=h$.
\end{proof}

\begin{lemma}[Monotone rounding]\label{lem:round}
For an integer $f$, let
\[
\eta_f(y)=\fl{\frac{y(10^4+f)+5000}{10^4}}.
\]
If $f\ge0$, then $\eta_f$ is nondecreasing on the integers and satisfies $\eta_f(y)\ge y$ for every integer $y\ge0$. Monotonicity also holds
for every $f\ge-10^4$.
\end{lemma}
\begin{proof}
For $f\ge-10^4$ the argument of the floor is nondecreasing in $y$, and the floor is
nondecreasing. For $f\ge0$ and $y\ge0$ we have $y(10^4+f)+5000\ge 10^4y$.
\end{proof}

\begin{lemma}[Lower bound correctness, L3]\label{lem:L3}
Assume \ref{as:nonneg}--\ref{as:extrap}. For every parcel $x$ and service $s$,
\[
h_s(x)\le\beta_s(x)\le c_s(x).
\]
In particular $h_s(x)=\infty$ implies $c_s(x)=\infty$. The first inequality uses $f\ge0$;
the second needs only $f\ge-10^4$.
\end{lemma}
\begin{proof}
If $s$ is ineligible, then $c_s=\infty$ and there is nothing to show. Otherwise apply Lemma~\ref{lem:base} with $t''=t_\gamma(x)\le\hat t_s(x)$;
this holds because tiers only take maxima. It gives $B_s\ge h_s(x)$; since $B_s<\infty$ for an
eligible service, $h_s(x)=\infty$ forces $c_s(x)=\infty$. All amounts are
nonnegative by \ref{as:nonneg}. In particular, per-weight and percentage-of-base amounts
are roundings of nonnegative numbers, so $\Sigma^\Fu,\Sigma^N\ge0$ and $Q\ge0$. By
Lemma~\ref{lem:round},
$c^0_s\ge\eta_f(B_s+\Sigma^\Fu)\ge\eta_f(h_s(x))=\beta_s(x)\ge h_s(x)$. Finally
$c_s=c^0_s+\rnd{c^0_sQ/100}\ge c^0_s$, because the rounding of a nonnegative number is
nonnegative.
\end{proof}

\subsection{Surcharge-class refinement (RQ+)}
Call a predicate \emph{geometric for $\gamma$} if it reads only parcel-level features,
or also $\mathit{DAS}$ when Assumption~\ref{as:das} holds for $\gamma$. Let $\Pi_\gamma$ be the
set of geometric predicates that occur in the tiers, ineligibility predicates or
surcharges of the services of $\gamma$. The class of $x$ is
$C(x)=\{\pi\in\Pi_\gamma:\pi(x)\}$.

\begin{definition}[Class key]\label{def:classkey}
For $s\in\gamma$, a local signature $\sigma$ with $z=\sigma_s$, an interval $e$ and a class $C$,
let $\hat e=\max\bigl(\{e\}\cup\{e_\gamma(m^{(j)}_s):\pi^{(j)}_s\in C\}\bigr)$. If $z=\bot$, if $\psi_s\in
C$, or if $h_s(z,\hat e)=\infty$, set $\chi_s=\infty$. Otherwise
\[
\chi_s(\sigma,e,C)=h_s(z,\hat e)+\sum_{(g,\varphi)}\max\bigl(\{0\}\cup\{A_r[z]: r\in\mathcal
R_s\text{ fixed},\ \pi_r\in C,\ (g_r,\varphi_r)=(g,\varphi)\}\bigr).
\]
Write $\chi_s=h_s(z,\hat e)+\chi^\Fu_s+\chi^N_s$, where $\chi^\Fu_s$ and $\chi^N_s$ sum over the
fuel and non-fuel keys. The row is \emph{non-fuel-flagged} (NF) if some service with
$\chi_s<\infty$ has $\chi^N_s>0$. The bound is $\beta^+_s=\chi_s$ for NF rows and
$\beta^+_s=\eta_f(\chi_s)$ otherwise.
\end{definition}

\begin{lemma}[L3 for RQ+]\label{lem:L3plus}
Assume \ref{as:nonneg}--\ref{as:extrap} and \ref{as:das}. For every $x$ and
$s\in\gamma$, $\beta^+_s(\sigma(x),e_x,C(x))\le c_s(x)$.
\end{lemma}
\begin{proof}
A geometric predicate has the same value in $C(x)$ and in the cost evaluation, because it
reads only features that do not depend on the service; for $\mathit{DAS}$ this uses
\ref{as:das}. If $\chi_s=\infty$ because of $z=\bot$ or $\psi_s\in C(x)$, then $s$ is
ineligible. Let $t''=\max\bigl(\{t_\gamma\}\cup\{m^{(j)}_s:\pi^{(j)}_s\in C(x)\}\bigr)$. Every tier in
$C(x)$ fires in the evaluation, so $t''\le\hat t_s(x)$. By Lemma~\ref{lem:band}(c),
$e_\gamma(t'')=\hat e$. If $s$ is eligible, Lemma~\ref{lem:base} gives $B_s\ge h_s(z,\hat e)$, and
in particular $h_s(z,\hat e)<\infty$. If $h_s(z,\hat e)=\infty$, the same lemma shows that $s$
is ineligible. For every key, the maximum over all firing rules of the key is at least
the maximum over its firing fixed geometric rules, since all amounts are nonnegative.
Hence $\Sigma^\Fu\ge\chi^\Fu_s$ and $\Sigma^N\ge\chi^N_s$. For an NF row,
$c_s\ge c^0_s\ge\eta_f(B_s+\Sigma^\Fu)+\Sigma^N\ge B_s+\Sigma^\Fu+\Sigma^N\ge\chi_s$. For a row that is
not NF, every service with finite key has $\chi^N_s=0$, so
$\eta_f(\chi_s)=\eta_f(h+\chi^\Fu_s)\le\eta_f(B_s+\Sigma^\Fu)\le c^0_s\le c_s$.
\end{proof}

\section{Query Exactness and Instance Optimality (T4)}\label{sec:T4}

\begin{definition}[Threshold algorithm TA]\label{def:TA}
On input $x$:
\begin{enumerate}[leftmargin=*]
\item Locate the elementary key interval of $\kappa$ and its signature $\sigma$. For every
group $\gamma$, compute $t_\gamma(x)$, $e_x=e_\gamma(t_\gamma)$, the fuel ratio $f_\gamma$ and the
list $\mathcal L_\gamma=\mathcal L_\gamma(\sigma|_\gamma,e_x)$. Set a pointer $p_\gamma$ to the first
entry, and drop the group if its first entry has $h=\infty$.
\item Set $(\mathit{best},\mathit{arg})=(\infty,-1)$. While some group remains: let $\gamma$ be
a remaining group whose head has the smallest $\beta$. If $\beta_{\mathrm{head}}>\mathit{best}$,
stop. Otherwise evaluate $v=c_s(x)$ for the head $s$, and replace
$(\mathit{best},\mathit{arg})$ by $(v,s)$ if $v<\mathit{best}$, or if $v=\mathit{best}$ and
$s<\mathit{arg}$. Advance $p_\gamma$, and drop $\gamma$ if its list is exhausted or its next
entry has $h=\infty$.
\item Return $(\mathit{best},\mathit{arg})$, where $(\infty,-1)$ encodes $\bot$.
\end{enumerate}
\end{definition}

\begin{definition}[Access model]\label{def:access}
An algorithm receives $x$ and the index. A \emph{sorted access} to a list returns the next
entry $(s,h_s)$, from which $\beta_s(x)$ follows, at cost $c_S$. An \emph{evaluation} of $s$
returns $c_s(x)$ at cost $c_R$. Nothing else about the costs is observable. The set of
instances consistent with the bounds is
$\mathcal I(\beta)=\{c'\in(\mathbb N\cup\{\infty\})^{\SC}: c'_s\ge\beta_s(x)\ \forall s\}$. An
algorithm is \emph{correct} if every run returns $\min_{\mathrm{lex}}$ of the finite pairs of
$c'$, or $\bot$, for every $c'\in\mathcal I(\beta)$. For randomized algorithms we require
correctness of every run (Las Vegas). This is the sorted/random access model of threshold
algorithms~\cite{fagin2003optimal}. As in minimal probing~\cite{chang2002minimal}, the
grade of an entry is replaced by a lower bound and a probe reveals the exact value.
\end{definition}

\begin{theorem}[Exactness and instance optimality, T4]\label{thm:T4}
Assume \ref{as:nonneg}--\ref{as:sid}, and let $(c^*,s^*)=\cfs(x)$.
\begin{enumerate}[label=(\alph*)]
\item TA returns $\cfs(x)$. It agrees with the brute force that scans the services in
identifier order and keeps a strictly smaller cost, including the tie rule and the answer
$\bot$.
\item If $\cfs(x)\ne\bot$, TA evaluates exactly the services in
$\{s:\beta_s(x)\le c^*\}$. If $\cfs(x)=\bot$, it evaluates exactly those with
$\beta_s(x)<\infty$. It performs at most (number of evaluations)$+|\Gamma|$ sorted accesses.
\item Let $M(x)=\{s:(\beta_s,s)\lexlt(c^*,s^*)\}\cup\{s^*\}$ if $\cfs(x)\ne\bot$, and
$M(x)=\{s:\beta_s<\infty\}$ otherwise. Every run of every correct algorithm on any
$c\in\mathcal I(\beta)$ with $\cfs=(c^*,s^*)$ evaluates every service of $M(x)$.
\item Let $\tau(x)=|\{s>s^*:\beta_s(x)=c^*\}|$, and let $\tau(x)=0$ if $\cfs(x)=\bot$. TA makes
exactly $|M(x)|+\tau(x)$ evaluations. For every correct algorithm $A$,
\[
\mathrm{cost}(\mathrm{TA},x)\le\Bigl(1+\frac{c_S}{c_R}\Bigr)\mathrm{cost}(A,x)+(c_R+c_S)\,\tau(x)+c_S|\Gamma|.
\]
Hence, on the inputs with $\tau(x)=0$, TA is instance optimal in the sense
of~\cite{fagin2003optimal} with optimality ratio $1+c_S/c_R$ and additive constant
$c_S|\Gamma|$, and its number of evaluations is the minimum over all correct algorithms. If
only evaluations are counted, the optimality ratio is $1$ with additive term $\tau(x)$.
\end{enumerate}
\end{theorem}
\begin{proof}
\emph{Order of pops.} Within a group, $\beta_s=\eta_f(h_s)$ with a common $f$
(Lemma~\ref{lem:L2}). By Lemma~\ref{lem:round}, every list is therefore nondecreasing in
$\beta$. TA always pops the smallest head, so the sequence of popped bounds is
nondecreasing, as in a merge of sorted lists.

\emph{Incumbent invariant.} After every step, $(\mathit{best},\mathit{arg})$ is the
lexicographic minimum of the finite pairs of the evaluated services, or $(\infty,-1)$ if
there is none. This holds initially. An infinite $v$ never replaces the incumbent: if
$\mathit{best}=\infty$, then $\mathit{arg}=-1$ and $s<-1$ fails. A finite $v$ replaces it
exactly when $(v,s)\lexlt(\mathit{best},\mathit{arg})$.

(a) Consider the moment TA stops or runs out of groups. Let $s$ be an unevaluated
service. Either $s$ lies after a dropped entry with $h=\infty$; then $h_s=\infty$ because the
lists are sorted, and $c_s=\infty$ by Lemma~\ref{lem:L3}. Or $s$ lies in a remaining list;
then $c_s\ge\beta_s\ge\beta_{\mathrm{head}}>\mathit{best}$. In both cases $(c_s,s)$ is not
lexicographically smaller than the incumbent, so by the invariant the incumbent is the
lexicographic minimum of all finite pairs, which is $\cfs(x)$. The brute force computes
the same minimum: it scans in identifier order by \ref{as:sid}, and a strict comparison
keeps the first, hence smallest, identifier among equal costs.

(b) Let $\cfs(x)=(c^*,s^*)$. When TA stops, $\mathit{best}=c^*$ by (a). Every service
with finite $h$ that was not evaluated lies in a remaining list and has
$\beta_s>\mathit{best}=c^*$. If the loop ended because no group remained, every service with
finite $h$ was evaluated. Hence every $s$ with $\beta_s\le c^*$ is evaluated. Conversely,
let $\beta_s>c^*$. Since $\beta_{s^*}\le c_{s^*}=c^*<\beta_s$ and pops are nondecreasing,
$s^*$ is popped and evaluated before $s$. From then on $\mathit{best}\le c^*$, and $\mathit{best}$
is the cost of an evaluated service, so $\mathit{best}\ge c^*$. When $s$ becomes the
smallest head, $\beta_s>\mathit{best}$ and TA stops. If $\cfs(x)=\bot$, then $\mathit{best}$
stays $\infty$, the stop test never fires, and exactly the services with finite $h$,
equivalently finite $\beta$, are evaluated. For the sorted accesses, each group reads its
first entry once, and each evaluation is followed by at most one further read.

(c) Suppose a run of a correct $A$ on $c$ never evaluates some $s\in M(x)$. Define $c'$ by
$c'_u=c_u$ for $u\ne s$, and set $c'_s=\beta_s$ if $s\ne s^*$, or $c'_{s^*}=c^*+1$ if $s=s^*$. If
$s\ne s^*$, then $\beta_s\le c^*<\infty$, or $\beta_s<\infty$ in the case $\cfs=\bot$, so $c'_s$ is
finite. Thus $c'\in\mathcal I(\beta)$. Every sorted access and every evaluation of the run
returns the same value under $c'$ as under $c$: the bounds are unchanged, and $s$ is never
evaluated. Hence the run returns the same answer, $(c^*,s^*)$ or $\bot$. But the correct
answer for $c'$ differs. If $s\ne s^*$ and $\cfs(x)\ne\bot$, then
$(c'_s,s)=(\beta_s,s)\lexlt(c^*,s^*)$, so the minimum for $c'$ is strictly smaller. If
$\cfs(x)=\bot$, then $c'$ has a finite pair. If $s=s^*$, every pair other than that of $s^*$
is unchanged and differs from $(c^*,s^*)$, and the pair of $s^*$ is now $(c^*+1,s^*)$; so no
pair equals $(c^*,s^*)$. This contradicts the correctness of $A$.

(d) We have $(\beta_s,s)\lexlt(c^*,s^*)$ iff $\beta_s<c^*$, or $\beta_s=c^*$ and $s<s^*$. Moreover,
$s^*$ belongs to the first set iff $\beta_{s^*}<c^*$. Hence
$|M|=|\{\beta_s<c^*\}|+|\{\beta_s=c^*,s\le s^*\}|$, while by (b) TA evaluates
$|\{\beta_s<c^*\}|+|\{\beta_s=c^*\}|=|M|+\tau$ services. Its cost is therefore at most
$c_R(|M|+\tau)+c_S(|M|+\tau+|\Gamma|)$. By (c), $\mathrm{cost}(A,x)\ge c_R|M|$, which gives
the inequality.
\end{proof}

\begin{remark}[Tie term]
The term $\tau(x)$ is unavoidable for any algorithm that pops services in bound order and
stops on a strict comparison. Sorting the lists by $(\beta,s)$ instead of $(h,s)$ does not
remove it in general, because $\eta_f$ may map distinct $h$ to the same $\beta$, and a
fuel-dependent order would contradict Lemma~\ref{lem:L2}. On instances where no other
service has a bound exactly equal to the optimal cost, $\tau=0$ and TA is exactly optimal
in evaluations.
\end{remark}

\begin{proposition}[RQ+ query]\label{prop:rqplus}
Assume \ref{as:nonneg}--\ref{as:sid} and \ref{as:das}. TA run on the RQ+ rows, whose lists are
sorted by $(\chi_s,s)$ with bound $\beta^+$, returns $\cfs(x)$. Statements (b)--(d) of
Theorem~\ref{thm:T4} hold with $\beta^+$ in place of $\beta$.
\end{proposition}
\begin{proof}
Within a row, $\beta^+$ is either $\chi$ or $\eta_f(\chi)$ with a common $f$, so the rows are
nondecreasing in $\beta^+$. Lemma~\ref{lem:L3plus} replaces Lemma~\ref{lem:L3}, and a
service with $\chi_s=\infty$ has $c_s=\infty$ by the proof of that lemma. The rest of the
proof of Theorem~\ref{thm:T4} uses only these two properties.
\end{proof}

\begin{proposition}[Batch query]\label{prop:batch}
Assume \ref{as:nonneg}--\ref{as:sid}. The vectorised algorithm, which in round $r$ evaluates
the rank-$r$ entry of every group for every parcel that is still active in that group and
deactivates a parcel in a group when the rank-$r$ bound is $\infty$ or exceeds the current
best, returns $\cfs(x)$ for every parcel.
\end{proposition}
\begin{proof}
The incumbent invariant holds because the lexicographic-minimum update is commutative
and associative. A group is deactivated for a parcel at rank $r$ only if
$\beta_r>\mathit{best}_{\mathrm{now}}\ge\mathit{best}_{\mathrm{final}}$ or $\beta_r=\infty$. All
later ranks have bounds at least $\beta_r$, so their costs exceed the final best or are
infinite. The argument of Theorem~\ref{thm:T4}(a) then applies.
\end{proof}

\begin{remark}[Truncated lists]
With list truncation to the first $\kappa$ entries, the implementation evaluates all
remaining services of a group when the truncated list is exhausted. Exactness is
preserved, because unlisted services have bounds at least as large as the last listed
bound. Statement (b) of Theorem~\ref{thm:T4} no longer holds. The surcharge-class index
requires untruncated lists (Appendix~\ref{app:mismatch}, item M11).
\end{remark}

\section{Index Updates}\label{sec:updates}

\begin{proposition}[Price edit]\label{prop:priceedit}
After $P_s[z,k]\leftarrow v$, the in-place repair recomputes $h_s(z,\cdot)$, and for every
local signature with $\sigma_s=z$ and every $e$ whose hull value changed, it removes $s$ and
reinserts it at the first position whose entry is lexicographically after
$(h^{\mathrm{new}},s)$. The resulting lists equal those of a fresh build. The repair takes
$O(|\{\sigma\in\Sigma_\gamma:\sigma_s=z\}|\cdot|E^\Delta|\cdot n_\gamma+K_s+m_\gamma\log K_s)$ time
with arrays, where $E^\Delta$ is the set of intervals whose hull value changed.
\end{proposition}
\begin{proof}
The value $h_{s'}(z',e)$ depends only on $(s',z',e)$ and on the prices of $s'$. Hence only
the entries of $s$ in lists with $\sigma_s=z$ can change, and only at intervals in
$E^\Delta$. Removal followed by insertion at the stated position restores the order by
$(h,s)$, which is strict because identifiers are distinct. The time bound counts the
hull computation and one array shift per touched list.
\end{proof}

For RQ+, the edit invalidates the materialised rows of the local signatures with
$\sigma_s=z$. These rows are rebuilt from the current data on their next access, so the
answers remain exact by Proposition~\ref{prop:rqplus}. A fuel edit changes only the
schedule, which takes $O(1)$ index work beyond the array insertion, by Lemma~\ref{lem:L2}.

\begin{proposition}[Zone edit]\label{prop:zoneedit}
After the edit $Z_s(\kappa)\leftarrow z'$ for $\kappa\in[l,u]$, the repair inserts the cuts $l$
and $u+1$ and, for every elementary interval inside $[l,u]$, replaces its signature by the
signature with $\sigma_s=z'$. Existing global and local signatures are reused, and new local
signatures receive freshly built lists. For every key, the resulting index returns the
same lists as a fresh build of the edited instance, so Theorem~\ref{thm:T4} continues to
hold.
\end{proposition}
\begin{proof}
Outside $[l,u]$ no zone map changes. The new cuts make every zone map constant on every
elementary interval (Proposition~\ref{prop:P1}). A list depends only on the local zone vector
and $e$, so reusing an existing local signature with the same vector, or building fresh
lists from current data, yields the lists of a fresh build. A fresh build may number the
signatures differently and merge adjacent intervals, which does not affect any answer.
\end{proof}

\section{Exact Maintenance (T7)}\label{sec:T7}

\begin{definition}[Top-2 state]\label{def:top2}
For a cost vector $c$ of a parcel, let $F(c)$ be the set of finite pairs. Then
$\mathrm{Top2}(c)=(p_1,p_2)$, where $p_1=\min_{\mathrm{lex}}F(c)$ and
$p_2=\min_{\mathrm{lex}}(F(c)\setminus\{p_1\})$; missing entries equal the padding
$\bot=(\infty,-1)$, which we order after every finite pair. We write $p_i=(c_i,s_i)$.
\end{definition}

\begin{definition}[Repair]\label{def:repair}
Given the state $(p_1,p_2)$ of a parcel, a service $s$ and its new cost $c'$, let
$p=(c',s)$:
\begin{enumerate}[label=(R\arabic*),leftmargin=*]
\item $s=s_1$: if $c'<\infty$ and $p\lexlt p_2$, set $p_1\leftarrow p$; otherwise
\emph{recompute}.
\item $s=s_2$: if $p\lexlt p_1$, set $(p_1,p_2)\leftarrow(p,p_1)$; else if $c'\le c_2$, set
$p_2\leftarrow p$; otherwise \emph{recompute}.
\item $s\notin\{s_1,s_2\}$ and $c'<\infty$: if $p\lexlt p_1$, set $(p_1,p_2)\leftarrow(p,p_1)$;
else if $p\lexlt p_2$, set $p_2\leftarrow p$; otherwise do nothing.
\item $s\notin\{s_1,s_2\}$ and $c'=\infty$: do nothing.
\end{enumerate}
To \emph{recompute} means to set the state to $\mathrm{Top2}$ of the current costs of all
services.
\end{definition}

\begin{lemma}[Single-service repair]\label{lem:repair1}
Let $(p_1,p_2)=\mathrm{Top2}(c)$, and let $c'$ differ from $c$ at most in coordinate $s$. Then
the repair with $c'_s$ yields $\mathrm{Top2}(c')$.
\end{lemma}
\begin{proof}
We use the fact (F): every finite pair of $c$ other than $p_1,p_2$ is lexicographically
after $p_2$. Pairs of services other than $s$ are the same under $c$ and $c'$.
(R1) The pair of $s$ was $p_1$. If $c'<\infty$ and $p\lexlt p_2$, then $p$ precedes $p_2$ and,
by (F), every other pair, so $\mathrm{Top2}(c')=(p,p_2)$. The other cases are recomputed.
(R2) The pair of $s$ was $p_2=(c_2,s)$. If $p\lexlt p_1$, then $p\lexlt p_1\lexlt$ every other
pair, so $\mathrm{Top2}(c')=(p,p_1)$. Otherwise $p_1\lexlt p$, since $p\ne p_1$. If moreover
$c'\le c_2$, then $c'$ is finite and $p\lexle(c_2,s)$, which precedes every other pair by
(F). Hence $\mathrm{Top2}(c')=(p_1,p)$. The remaining case is recomputed.
(R3) The old pair of $s$, if finite, came after $p_2$. If $p\lexlt p_1$, then
$\mathrm{Top2}(c')=(p,p_1)$. If $p_1\lexlt p\lexlt p_2$, then it is $(p_1,p)$. Otherwise $p$
comes after $p_2$, and both the old and the new pair of $s$ come after $p_2$, so the state
is unchanged.
(R4) The pair of $s$ leaves $F$. It came after $p_2$, so the state is unchanged.
\end{proof}

\begin{lemma}[Several edited services]\label{lem:compose}
Suppose that an edit changes the costs of $S_\Delta=\{u_1,\dots,u_r\}$ from $c$ to $c'$, and let all repairs and recomputations use $c'$. If each parcel is repaired once for
each $u\in S_\Delta$, in any order, its final state is $\mathrm{Top2}(c')$. The same holds with
\emph{deferred recomputation}: a parcel that reaches a recompute branch is only marked, and
the remaining repairs for it are applied to its current state (or skipped). After all
services of $S_\Delta$ have been processed, every marked parcel is recomputed once. This
recomputation may reuse the already evaluated costs $c'_u(x)$ for $u\in S_\Delta$.
\end{lemma}
\begin{proof}
We prove the invariant that after the repairs for $u_1,\dots,u_i$, the state equals
$\mathrm{Top2}(c^{(i)})$ for a vector $c^{(i)}$ with $c^{(i)}_u=c'_u$ for $u\notin S_\Delta$ and for
$u\in\{u_1,\dots,u_i\}$. For $i=0$, take $c^{(0)}=c$, which agrees with $c'$ outside $S_\Delta$.
For the step, if the repair for $u_{i+1}$ does not recompute, Lemma~\ref{lem:repair1}
applied to $c^{(i)}$ and $c^{(i)}[u_{i+1}\leftarrow c'_{u_{i+1}}]$ yields the invariant. If it
recomputes, the state becomes $\mathrm{Top2}(c')$, and $c^{(i+1)}=c'$ satisfies the
invariant. For $i=r$ we have $c^{(r)}=c'$.

For deferred recomputation, an unmarked parcel has gone only through non-recompute
repairs, so the invariant above holds for it at every step, and its final state is
$\mathrm{Top2}(c')$. A marked parcel is finally set to $\mathrm{Top2}(c')$ by a recomputation
over all services with the edited instance, which does not read the intermediate state. A
recomputation that takes $c'_u(x)$ for $u\in S_\Delta$ from the repair phase and evaluates the
other services computes the same top-2, because the top-2 of a set of pairs does not
depend on the order in which the pairs are inserted.
\end{proof}

\begin{definition}[Supports]\label{def:supports}
Let $X$ be the historical parcels, $t_\gamma(x)$ their pre-tier weights and
$e_x=e_\gamma(t_\gamma(x))$. A \emph{bucket} $(\sigma,e)$ is the set of parcels whose signature
is $\sigma$ and whose interval in the group of the edited service is $e$.
\begin{enumerate}[label=(S\arabic*),leftmargin=*]
\item \emph{Cell $(s,z,k)$.} Let $\bar b=b_{s,k}$, except $\bar b=+\infty$ if $s$
extrapolates and $k=K_s$. Let $I_k=\{e\in\{1,\dots,m+1\}:b_{s,k-1}<E_e\le\bar b\}$. Then
\[
U_{\mathrm{cell}}=\{x:Z_s(\kappa)=z,\ e_x\in I_k\}\cup\bigcup_{j:\ b_{s,k-1}<m^{(j)}_s\le\bar b}\{x:Z_s(\kappa)=z,\ \pi^{(j)}_s(x)\}.
\]
\item \emph{Band $(s,k)$ over zones $Z'$, or any set of cells of one service.} The union of
the cell supports.
\item \emph{Table rewrite of $s$, such as a general rate increase.}
$U_{\mathrm{srv}}=\{x:Z_s(\kappa)\ne\bot\}$.
\item \emph{Rule edit of $s$}, which changes $\mathcal R_s$ or $\psi_s$ but not zones, prices,
the group or the tiers. If $\psi_s$ changes, or if a changed rule is a
percentage-of-total rule, the support is $U_{\mathrm{srv}}$. Otherwise, let $\Delta\mathcal R$ be
the symmetric multiset difference of the old and the new rule lists, comparing all fields
of a rule. Then
$U_{\mathrm{rule}}=\{x\in U_{\mathrm{srv}}:\exists r\in\Delta\mathcal R,\ \pi_r(x)\}$, where $\pi_r$
is evaluated as in Definition~\ref{def:cost}.
\item \emph{Zone edit} $Z_s(\kappa)\leftarrow z'$ for $\kappa\in[l,u]$:
$U_{\mathrm{zone}}=\{x:\kappa\in[l,u],\ Z^{\mathrm{old}}_s(\kappa)\ne z'\}$.
\item \emph{Fuel edit} of class $\phi$, segment $w$: let $D_w=[\sigma_{\phi,w},\sigma_{\phi,w+1})$,
with $\sigma_{\phi,0}$ replaced by $-\infty$ for $w=0$ and $\sigma_{\phi,W}=+\infty$. Then
$U_{\mathrm{fuel}}=\{x:d(x)\in D_w\}$, and $S_\Delta=\{u:\phi_u=\phi\}$. Inserting a new start
$\sigma'$ into segment $w$ gives the support $\{x:d(x)\in[\sigma',\sigma_{\phi,w+1})\}$.
\item \emph{Service insertion or deletion:} $U_{\mathrm{srv}}$, with the cost of the inserted
service before the edit, or of the deleted service after it, taken as $\infty$.
\end{enumerate}
\end{definition}

\begin{lemma}[Cell support]\label{lem:cellsupport}
Assume \ref{as:extrap}, \ref{as:tier} and \ref{as:spacing}. If $c_s(x)$ changes when only
$P_s[z,k]$ changes, then $x\in U_{\mathrm{cell}}$.
\end{lemma}
\begin{proof}
The cost reads the price table only through $B_s$, which is also the only input of
percentage-of-base amounts. $B_s$ reads $P_s[z,k]$ only if $Z_s(\kappa)=z$ and either (i)
$k_s(\hat t)=k$ and $s$ does not extrapolate at $x$, or (ii) $s$ extrapolates with
$\hat t>\lambda_s$ and $k=K_s$. Let $t=t_\gamma(x)$. By \ref{as:tier}, firing of tiers does not
depend on $T$, so $\hat t=\max(\{t\}\cup\{m^{(j)}_s:\pi^{(j)}_s(x)\})$.

\emph{Case $\hat t=t$.} In case (i), if some breakpoint is at least $t-\epsilon$,
Lemma~\ref{lem:band}(a) gives $b_{s,k-1}<E_{e_x}\le b_{s,k}\le\bar b$, so $e_x\in I_k$.
Otherwise $k_s(t)=K_s=k$ and $t-\epsilon>b_{s,K_s}$ by Lemma~\ref{lem:band}(d). Then a
non-extrapolating $s$ is ineligible both before and after the edit, so its cost cannot
change. If $s$ extrapolates, then $\bar b=\infty$ and $E_{e_x}\ge t-\epsilon>b_{s,K_s}-\epsilon\ge
b_{s,K_s-1}$; for $e_x=m+1$ this holds trivially, and otherwise it uses \ref{as:spacing}. So
$e_x\in I_{K_s}$. In case (ii), $t>\lambda_s\ge b_{s,K_s}$ by \ref{as:extrap}, and the same
argument gives $e_x\in I_{K_s}$.

\emph{Case $\hat t>t$.} Then $\hat t=m^{(j)}_s$ for a firing tier $j$. In case (i),
$k_s(m^{(j)}_s)=k$ means $b_{s,k-1}<m^{(j)}_s-\epsilon$ and, if $k<K_s$ or no clipping occurs,
$m^{(j)}_s-\epsilon\le b_{s,k}$. Hence $b_{s,k-1}<m^{(j)}_s$, and $m^{(j)}_s\le b_{s,k}+\epsilon$,
which together with \ref{as:tier} gives $m^{(j)}_s\le b_{s,k}$. If clipping occurs, then
$k=K_s$, and either $s$ extrapolates and $\bar b=\infty$, or $s$ is ineligible both before
and after the edit. In case (ii), $m^{(j)}_s>\lambda_s\ge b_{s,K_s}>b_{s,K_s-1}$ and
$\bar b=\infty$. In every case $b_{s,k-1}<m^{(j)}_s\le\bar b$, and $x$ belongs to the tier
term of $U_{\mathrm{cell}}$, because by \ref{as:tier} the tier predicate evaluated without
service-dependent features agrees with the evaluation.
\end{proof}

\begin{lemma}[Other supports]\label{lem:othersupports}
If the costs change on parcel $x$ under an edit of type (S2)--(S7), then $x$ lies in the
corresponding support.
\end{lemma}
\begin{proof}
(S2) A set of cells of one service changes $c_s(x)$ only if one of the cells is read at
$x$; apply Lemma~\ref{lem:cellsupport} to each cell. (S3) If $Z_s(\kappa)=\bot$, then
$c_s(x)=\infty$ before and after the edit. (S4) The zones, $\hat t$ and $B_s$ are unchanged.
If $x$ is in no support, then no rule of $\Delta\mathcal R$ fires at $x$. The firing rules at $x$,
with their amounts and keys, are then the same before and after the edit, because the
unchanged rules are identical and the changed ones do not fire. Since every
$\mu_{g,\varphi}$ is a maximum over firing rules and $\psi_s$ is unchanged,
(\ref{eq:cost}) is unchanged. (S5) $c_s$ reads $\kappa$ only through $Z_s(\kappa)$, which also
supplies the $\mathit{ZONE}$ feature, and through $\mathit{DAS}_s(\kappa)$, which is
unchanged. (S6) $f_\phi(d)$ changes exactly for $d\in D_w$, and the edit changes no other
input of any service. (S7) A service with $Z_s=\bot$ has cost $\infty$ before and after.
\end{proof}

\begin{theorem}[Exact maintenance, T7]\label{thm:T7}
Let an edit $\Delta$ change the services in $S_\Delta$, and let
$D_\Delta=\{x\in X:\exists u\in S_\Delta,\ c_u(x)\ne c'_u(x)\}$. For every $U$ with
$D_\Delta\subseteq U\subseteq X$, repairing every $x\in U$ once for each $u\in S_\Delta$, in any
order and with the edited instance, makes the state of every $x\in X$ equal to
$\mathrm{Top2}(c'(x))$. Under \ref{as:extrap}, \ref{as:tier} and \ref{as:spacing}, the
supports of Definition~\ref{def:supports} contain $D_\Delta$. The repair evaluates
$|S_\Delta|\cdot|U|$ costs and recomputes the set $R\subseteq U$ of parcels that reach a
recompute branch. With deferred recomputation (Lemma~\ref{lem:compose}), each parcel of
$R$ is recomputed once per edit, at a cost of $|\SC|-|S_\Delta|$ further evaluations. A parcel reaches a recompute branch only if the edited service was its best and became
infeasible or lexicographically worse than the runner-up (R1), or was its runner-up and
became more expensive than before without overtaking the best (R2).
\end{theorem}
\begin{proof}
For $x\notin U$, all costs are unchanged, so the old state equals $\mathrm{Top2}(c'(x))$. For
$x\in U$, apply Lemma~\ref{lem:compose}. The containment of the supports is
Lemmas~\ref{lem:cellsupport} and~\ref{lem:othersupports}. The cost statement counts one
evaluation per parcel and service, plus one deferred recomputation per element of $R$. The
characterisation of $R$ follows from (R1) and (R2).
\end{proof}

\begin{remark}[Granularity]
$U_{\mathrm{cell}}$ is a union of whole buckets plus tier sets. A parcel in a selected bucket whose tier moves it out of band $k$, or for which $s$ is
ineligible, is included although its cost does not change; its repair is a no-op.
Exactness never depends on minimality.
\end{remark}

\section{Complexity}\label{sec:complexity}

Let $n=|\SC|$, $N\le 1+\sum_sR_s$ the number of elementary key intervals, $\Gamma$ the set of
groups, $n_\gamma$ the size of group $\gamma$, $m_\gamma=|E_\gamma|\le\sum_{s\in\gamma}K_s$,
$\Sigma_\gamma$ the local signatures, and $Z^\gamma_s$ the number of distinct zones of $s$ that
occur in $\Sigma_\gamma$. Let $C_{\mathrm{ev}}(s)$ be the cost of one evaluation, which is
$O(\log R_s+\log K_s+\log|\mathit{DAS}_s|+\sum_{r}|\pi_r|+|\pi^{(\cdot)}_s|+|\psi_s|)$, where $|\pi|$
is the number of conditions of a predicate.

\begin{proposition}[Build, space, query and maintenance cost]\label{prop:complexity}
\begin{enumerate}[label=(\alph*)]
\item \emph{Build:} the index is built in time
\[
O\Bigl(nN\log N+\sum_\gamma\Bigl[\sum_{s\in\gamma}Z^\gamma_s(K_s+m_\gamma\log K_s)+|\Sigma_\gamma|(m_\gamma+1)\,n_\gamma\log n_\gamma\Bigr]\Bigr).
\]
\item \emph{Space:} $O\bigl(N+n|\Sigma|+\sum_\gamma|\Sigma_\gamma|(m_\gamma+1)\,n_\gamma\bigr)$ words. RQ+ adds
$O(n_\gamma)$ words per materialised row $(\sigma,e,C)$, and the number of such rows never
exceeds the number of distinct triples $(\sigma|_\gamma,e_x,C(x))$ among the queried parcels.
\item \emph{Query (TA):} $O\bigl(\log N+\sum_\gamma(C_\tau+\log m_\gamma)+N_{\mathrm{TA}}(x)\,(|\Gamma|+C_{\mathrm{ev}})\bigr)$,
where $C_\tau$ is the cost of computing $t_\gamma$ and $N_{\mathrm{TA}}(x)=|\{s:\beta_s(x)\le c^*\}|$ by
Theorem~\ref{thm:T4}(b). With a heap over the group heads, the factor $|\Gamma|$ becomes
$\log|\Gamma|$.
\item \emph{Maintenance:} initialisation takes $O(|X|\,\sum_sC_{\mathrm{ev}}(s)+|\Gamma|\,|X|\log|X|)$.
Enumerating $U_{\mathrm{cell}}$ takes $O(|\{\sigma:\sigma_s=z\}|\cdot|I_k|+|U|\log|U|)$ given the
cached tier sets. The repair costs $|U|$ evaluations of $s$ plus $|R|\cdot n$ evaluations.
\end{enumerate}
\end{proposition}
\begin{proof}
(a) Computing the zone vectors takes $n$ binary searches per elementary interval, and
deduplication is a sort of $N$ vectors of length $n$. For every service and every distinct
zone, the suffix minimum takes $O(K_s)$, and locating the $m_\gamma$ intervals takes
$O(m_\gamma\log K_s)$. Sorting takes $O(n_\gamma\log n_\gamma)$ per list.
(b) Count the stored arrays.
(c) Point location and the per-group set-up are followed, for each evaluated service, by
one selection among at most $|\Gamma|$ heads and one evaluation; the number of pops is
$N_{\mathrm{TA}}$.
(d) The initial top-2 evaluates every service on every parcel, and the buckets are sorted
once per group. Enumeration performs one dictionary lookup per candidate bucket and
deduplicates the result.
\end{proof}

\section{Envelope Complexity and General Predicates (T5, T6)}\label{sec:T5T6}

Throughout this section, fix a group $\gamma$, a local signature, the fuel ratio and the
values of all predicates. Assume that all surcharges are fixed amounts and that no
service extrapolates. Then $c_s(x)$ depends on the weight only through $k_s(\hat t_s)$.
Define the \emph{restricted fragment} by the additional condition that no tier fires,
so that $\hat t_s=t_\gamma$.

\begin{proposition}[One-dimensional envelope, T5 part 1]\label{prop:T5a}
In the restricted fragment, the map $t\mapsto\cfs$ restricted to $\gamma$ is constant on each
elementary interval $\{t:e_\gamma(t)=e\}$, so it has at most $m_\gamma+1\le\sum_{s\in\gamma}K_s+1$
pieces. The bound $m_\gamma+1$ is attained for every $m_\gamma$.
\end{proposition}
\begin{proof}
By Lemma~\ref{lem:band}(a), $k_s(t)$ is determined by $e_\gamma(t)$ for every $s\in\gamma$, so all
costs are constant on each elementary interval. For tightness, take services
$s_1,\dots,s_m$ with $s_1$ having the single breakpoint $1$ and price $1$, and $s_i$ for
$i\ge2$ having breakpoints $i-1,i$ and prices $3,1$. On $(i-1,i]$ the service $s_i$ costs
$1$, every $s_{i'}$ with $i'>i$ costs $3$, and every $s_{i'}$ with $i'<i$ is ineligible. Hence
the argmin is $s_i$ on $(i-1,i]$, and $\bot$ on $(m,\infty)$. This gives $m+1$ pieces with
$E_\gamma=\{1,\dots,m\}$.
\end{proof}

\begin{proposition}[Two-dimensional overlay, T5 part 2]\label{prop:T5b}
For every rounding rule, $t_\gamma(a,V)=\max(\alpha_\gamma(a),\nu_\gamma(V))$ with nondecreasing
$\alpha_\gamma,\nu_\gamma$. Hence, across all groups in the restricted fragment, the plane
$(a,V)$ is partitioned into at most $O(\Lambda^2)$ cells on which every cost is constant,
where $\Lambda=\sum_\gamma(m_\gamma+1)$. There are instances with two groups where the overlay
has $\Omega(\Lambda^2)$ cells.
\end{proposition}
\begin{proof}
For the carrier rule, $\alpha(a)=a$ and $\nu(V)=[V>\theta]V/D$. For the erp-style rule, the
rounding of $a$ and of $V$ are applied separately and then maximised. For the enterprise
rule in mode $1$, $r(\max(a,v))=\max(r(a),r(v))$ because $r$ is nondecreasing; modes $2$
and $3$ are the special cases with a constant $\nu$ or $\alpha$. Hence
$\{t_\gamma\le w\}=\{\alpha(a)\le w\}\cap\{\nu(V)\le w\}$ is a down-closed axis-parallel
rectangle, and the boundary between consecutive elementary intervals consists of at
most one vertical and one horizontal segment or ray. The overlay of all groups is
therefore an arrangement of at most $2\Lambda$ axis-parallel segments. Such an arrangement
has $O(\Lambda^2)$ vertices and hence, by Euler's formula, $O(\Lambda^2)$
faces~\cite{deberg2008computational}. Every cost is constant on each face.

For the lower bound, take two carrier groups with $\theta=0$, divisors $D_1>D_2$ with
$D_1/D_2>2K$, and $E_1=E_2=\{1,\dots,K\}$. In the wedge $V/D_1<a<V/D_2$ we have
$t_1=a$ and $t_2=V/D_2$. For $1\le i<j\le K$, the point $(i-\tfrac12,(j-\tfrac12)D_2)$ lies in
the wedge, because $(j-\tfrac12)D_2/D_1<1/2\le i-\tfrac12<j-\tfrac12$. Its pair of intervals is
$(i,j)$. Hence at least $K(K-1)/2=\Omega(\Lambda^2)$ distinct cells occur.
\end{proof}

\begin{conjecture}[T5 lower bound for the argmin]\label{conj:T5}
There are instances in the restricted fragment where the number of maximal connected
regions of the $(a,V)$ plane with a constant argmin is $\Omega(\Lambda^2)$.
\end{conjecture}
Proposition~\ref{prop:T5b} shows this lower bound for the overlay of cells, not for the
argmin. With one service per group, the argmin regions are separated by a threshold on
the two price sequences, and a checkerboard pattern of the argmin is impossible. A
construction with $\Theta(K)$ services per group is a plausible route.

\begin{proposition}[Class counts, T6 part 1]\label{prop:T6a}
(a) If the predicates of $\gamma$ are conjunctions and disjunctions of threshold conditions
on the parcel-level features, with $h_\ell$ distinct thresholds on feature $\ell$ and $g_c$
distinct values that matter for categorical feature $c$, then
$|\{C(x)\}|\le\prod_\ell(h_\ell+1)\prod_cg_c$.
(b) If the predicates are $H$ halfspaces in $\mathbb R^d$ in general position, then the
number of classes is at least $\sum_{i=0}^{d}\binom Hi=\Theta(H^d)$ for fixed $d$. Every index that
materialises one row per class eagerly therefore needs $\Omega(H^d)$ rows.
\end{proposition}
\begin{proof}
(a) A class is determined by the cell of the grid of thresholds and by the categorical
values, because every condition is constant on each grid cell. (b) Distinct full-dimensional
cells of the arrangement have distinct sign vectors, and their number is
$\sum_{i\le d}\binom Hi$~\cite{buck1943partition,zaslavsky1975facing}.
\end{proof}

\begin{conjecture}[T6]\label{conj:T6}
For general linear predicates in $d$ dimensions, there is no exact CFS index of size
polynomial in $H$ and $n$, independent of $d$, that answers queries with
$o(n)$ evaluations in the worst case.
\end{conjecture}
Proposition~\ref{prop:T6a}(b) excludes only eager per-class materialisation, which is the
design that RQ+ avoids by building rows lazily. A general lower bound would require a
cell-probe argument, which we do not claim.

\appendix
\section{Correspondence with the Code}\label{app:code}

Line numbers refer to the code as of 2026-11-10, after the consistency fixes, in
\code{code/rateq/} unless stated otherwise.

\begin{longtable}{>{\raggedright\arraybackslash}p{0.34\textwidth}>{\raggedright\arraybackslash}p{0.58\textwidth}}
\toprule
Theory & Implementation \\
\midrule
\endhead
Zone map $Z_s$ (Def.~\ref{def:program}(i)) & \code{model.py:157-160} \code{zone\_of}; scalar \code{index.py:78-83} \\
Billable weight $\tau_\gamma$ & \code{model.py:222-237} \code{billable}; \code{index.py:58-73} \code{billable1} \\
Group key $\gamma(s)$ & \code{model.py:153-155} \code{group\_key} \\
Tiers, lower first & \code{model.py:148-151} \code{min\_tiers}, \code{model.py:249-251}, \code{index.py:92-95} \\
Band $k_s$ with $\epsilon=10^{-9}$ & \code{model.py:261-262}, \code{index.py:103-105} \\
Extrapolation with cap & \code{model.py:255-269}, \code{index.py:99-112} \\
Eligibility & \code{model.py:252-254, 260, 270}, \code{index.py:78-114} \\
Surcharges, keys, maxima & \code{model.py:275-294}; \code{index.py:115-137} \\
Cost (\ref{eq:cost}) & \code{model.py:295-299}; \code{index.py:138-144} \\
CFS, brute force & \code{model.py:305-315}, \code{index.py:159-165} \\
Assumptions A1, A3 (build-time check) & \code{index.py:176-193} \\
Assumptions A5, A6 (maintenance check) & \code{maintain.py:54-68} \\
Assumption A7 (\code{das\_ok}) & \code{index.py:510-516} \\
Key partition, signatures & \code{index.py:203-214} \\
$E_\gamma$, hull $h$, lists & \code{index.py:221, 228-243, 245-265} \\
TA (Def.~\ref{def:TA}) & \code{index.py:279-320}; stop test \code{index.py:298} \\
Batch (Prop.~\ref{prop:batch}) & \code{index.py:322-379} \\
Price and zone edits & \code{index.py:382-414, 428-472} \\
Class key $\chi$, NF flag & \code{index.py:581-615}; geometric features \code{index.py:483-493} \\
RQ+ TA & \code{index.py:631-666}, \code{index.py:668-725} \\
Top-2, repair (R1)--(R4) & \code{maintain.py:21-40, 151-202} \\
Deferred recomputation (Lemma~\ref{lem:compose}) & \code{maintain.py:155-202} \code{apply\_many} \\
$U_{\mathrm{cell}}$ (S1), signature refresh & \code{maintain.py:70-87, 97-122}; \code{index.py:198, 471} \\
$U_{\mathrm{zone}}$ (S5) & \code{maintain.py:124-130} \\
$U_{\mathrm{srv}}$ (S3), $U_{\mathrm{fuel}}$ (S6) & \code{maintain.py:132-134, 136-149} \\
$U_{\mathrm{rule}}$ (S4) & \code{code/erp\_pilot.py:187-201, 213-227} \\
\bottomrule
\end{longtable}

\paragraph{Branch equivalence of the repair.} The code encodes $\bot$ as $(\infty,-1)$, and its
comparison treats $(\infty,-1)$ as smaller than $(\infty,s)$ for $s\ge0$. Every branch
nevertheless coincides with Definition~\ref{def:repair}. For $s=s_1$, the code keeps $p$ iff
$\neg((c_2,s_2)\lexlt(c',s))$. If $c'=\infty$, this is false whether or not $p_2$ is the
padding, so the code recomputes, as in (R1). For $s=s_2$, the promotion test
$(c',s)\lexlt(c_1,s_1)$ is false for $c'=\infty$, because $c_1$ is finite, and
$c'\le c_2$ is false as well, so the code recomputes. For other services with $c'=\infty$,
both tests are false, as in (R4).

\paragraph{Assumption check.} A scan of the compiled instances found no violation of
\ref{as:nonneg}--\ref{as:das}. The instances are the public cards, their scaled variants,
the synthetic and erp-style generators, and the enterprise US instance. In detail:
\begin{itemize}[leftmargin=*]
\item no price, amount, coefficient or fuel ratio is negative;
\item no extrapolating service has $\lambda_s<b_{s,K_s}$;
\item the identifier of every service equals its position;
\item no tier predicate reads $T$, $\mathit{ZONE}$, $\mathit{ZIP}$ or $\mathit{DAS}$;
\item no tier minimum lies within $\epsilon$ above a breakpoint;
\item every group that shares a delivery-area start array also shares the level array.
\end{itemize}
Since 2026-11-10 the code enforces these assumptions. \ref{as:nonneg} and \ref{as:extrap}
are asserted when the index is built, and \ref{as:tier} and \ref{as:spacing} when the
maintainer is constructed. \ref{as:das} is part of the \code{das\_ok} test.

\section{Code--Theory Discrepancies and Their Resolution}\label{app:mismatch}

The first version of this document (2026-11-09) listed twelve discrepancies. None of them
invalidated a reported result: each was either a latent precondition that held on all
instances, or a superset that affected only efficiency. The code was revised on
2026-11-10. The status below uses the revised line numbers. \emph{Resolved} means that
the code now implements the theory exactly. \emph{Enforced} means that the precondition is
asserted, so a violating instance is rejected. \emph{By design} means that the behaviour is
correct and accounted for in the theory.

\begingroup\raggedright
\begin{enumerate}[label=M\arabic*,leftmargin=*]
\item \textbf{Stop rule and ties}, by design. The locations are \code{index.py:298} and
\code{index.py:651}. The strict test $\beta_{\mathrm{head}}>\mathit{best}$ evaluates all services
with bound equal to $c^*$. This yields the additive term $\tau(x)$ of
Theorem~\ref{thm:T4}(d). The research plan's claim that every certificate-producing
algorithm evaluates $|\{s:\beta_s\le c^*\}|$ services is corrected: the lower bound is
$|M(x)|$, which is smaller by $\tau(x)$.
\item \textbf{Extra interval in $U_{\mathrm{cell}}$}, resolved (\code{maintain.py:101-108}). The
code now enumerates exactly $I_k=\{e:b_{s,k-1}<E_e\le\bar b\}$. Previously it also
enumerated the interval directly after $b_{s,k}$.
\item \textbf{Tier band test without $\epsilon$}, enforced (\code{maintain.py:117}). The test
is equivalent to $\epsilon$-band membership under the second part of \ref{as:tier}, which
is now asserted in \code{maintain.py:54-68}.
\item \textbf{Tier masks}, resolved and enforced (\code{maintain.py:54-68, 89-95}). The first
part of \ref{as:tier} is asserted, and the masks are evaluated on parcel-level features
only. Under that assumption, caching them by predicate identity across services is
exact.
\item \textbf{Stale signatures after zone edits}, resolved. \code{RQIndex.update\_zone}
increments a version counter (\code{index.py:198, 471}). \code{support\_cell} rebuilds the
parcel signatures and buckets when the version has changed
(\code{maintain.py:70-87, 97-100}). \code{support\_zone} (\code{maintain.py:124-130}) implements
$U_{\mathrm{zone}}$ of (S5) and must be called before the index edit.
\item \textbf{Fuel support at the ends of the schedule}, resolved. \code{support\_fuel}
(\code{maintain.py:136-149}) implements $D_w$ of (S6): segment $0$ extends to $-\infty$ and
the last segment to $+\infty$. The pilot calls it in \code{code/pilot.py:255-262}.
\item \textbf{Rule-level support}, resolved and enforced. The tiers are now applied in
tier order (\code{code/erp\_pilot.py:222-223}). The rule identity includes the fuel flag
(\code{code/erp\_pilot.py:187}). A replayed event is asserted to be a rule edit in the sense
of (S4), with zones, prices, group and tiers unchanged (\code{code/erp\_pilot.py:190-201}).
\item \textbf{Negative amounts}, enforced (\code{index.py:176-193}). \ref{as:nonneg} is
asserted at build time. Under it, the scalar evaluator (\code{index.py:131}) and the
vector evaluator (\code{model.py:289}) coincide.
\item \textbf{Extrapolation limit}, enforced (\code{index.py:176-193}). \ref{as:extrap} is
asserted at build time.
\item \textbf{Shared delivery areas}, resolved (\code{index.py:510-516}). The flag
\code{das\_ok} now requires the same start arrays and equal level arrays
(\ref{as:das}).
\item \textbf{RQ+ with truncation}, by design (\code{index.py:526-532}). The unsorted hull
is recovered from the lists, so RQ+ requires \code{kappa=None}, which is the default.
\item \textbf{Documentation}, resolved (\code{model.py:1-20} and the docstring of
\code{RQIndexPlus}, \code{index.py:495}). The documentation now matches the theory: no
monotonicity is assumed, and the full list of geometric features is given.
\end{enumerate}
\endgroup

\paragraph{Merged fuel-class update.} \code{Maintainer.apply\_many}
(\code{maintain.py:155-202}) implements the deferred recomputation of
Lemma~\ref{lem:compose}. It applies the repairs (R1)--(R4) once per service of the edited
class on a common support and marks the parcels that reach a recompute branch. It then
recomputes each marked parcel once and reuses the new costs of the edited services. The
single-service \code{apply} is \code{apply\_many} with one service.

\begin{thebibliography}{9}
\bibitem{fagin2003optimal}
R.~Fagin, A.~Lotem, and M.~Naor, ``Optimal aggregation algorithms for middleware,''
\emph{J. Comput. Syst. Sci.}, vol.~66, no.~4, pp.~614--656, 2003.
\bibitem{chang2002minimal}
K.~C.-C. Chang and S.-w. Hwang, ``Minimal probing: Supporting expensive predicates for
top-k queries,'' in \emph{Proc. ACM SIGMOD}, 2002, pp.~346--357.
\bibitem{buck1943partition}
R.~C. Buck, ``Partition of space,'' \emph{Amer. Math. Monthly}, vol.~50, no.~9,
pp.~541--544, 1943.
\bibitem{zaslavsky1975facing}
T.~Zaslavsky, ``Facing up to arrangements: Face-count formulas for partitions of space
by hyperplanes,'' \emph{Mem. Amer. Math. Soc.}, vol.~1, no.~154, 1975.
\bibitem{deberg2008computational}
M.~de~Berg, O.~Cheong, M.~van~Kreveld, and M.~Overmars, \emph{Computational Geometry:
Algorithms and Applications}, 3rd~ed. Springer, 2008.
\end{thebibliography}

\end{document}
